% from ega4-17-fr.tex
\oldpage[IV-4]{83}

Il reste donc à prouver que c) entraîne a). On peut se borner au cas où $S=\Spec(A)$, $Y=\Spec(B)$, $X=\Spec(C)$ sont affines. L'hypothèse implique (17.2.3, (i)) que $(\Omega^1_{X/S})_x=\Omega^1_{C_{\mathfrak m}/A_{\mathfrak s}}$ est un $C_{\mathfrak m}$-module libre: en effet (16.10.6) il existe des éléments $t_i$ ($1\leq i\leq r$) de $B_{\mathfrak p}=\sh{O}_{Y,y}$ tels que les différentielles $d_{B_{\mathfrak p}/A_{\mathfrak s}}(t_i)$ engendrent le $B_{\mathfrak p}$-module $\Omega^1_{B_{\mathfrak p}/A_{\mathfrak s}}$ et que leurs images dans $\Omega^1_{C_{\mathfrak m}/A_{\mathfrak s}}$ fassent partie d'une base de ce $C_{\mathfrak m}$-module libre. Comme $\Omega^1_{Y/S}$ (resp. $\Omega^1_{X/S}$) est un $\sh{O}_Y$-Module (resp. un $\sh{O}_X$-Module) de présentation finie (16.4.22), on peut, en remplaçant au besoin $X$ et $Y$ par des voisinages ouverts affines convenables de $x$ et $y$ respectivement, supposer que $\Omega^1_{C/A}$ est un $C$-module libre et que les $t_i$ sont les images d'éléments $s_i$ ($1\leq i\leq r$) de $B$ tels que les $d_{B/A}(s_i)$ engendrent le $B$-module $\Omega^1_{B/A}$ et que leurs images dans $\Omega^1_{C/A}$ fassent partie d'une base de ce $C$-module (Bourbaki, \emph{Alg. comm.}, chap. II, § 5, n\textsuperscript{o} 1, prop. 2). Soit $\varphi$ le $A$-homomorphisme de $B'=A[T_1,\ldots,T_r]$ dans $B$ tel que $\varphi(T_i)=s_i$ pour tout $i$; le di-homomorphisme correspondant $\Omega^1_{B'/A}\to\Omega^1_{B/A}$ (0, 20.5.2) transforme les $d_{B'/A}(T_i)$, qui forment une base de $\Omega^1_{B'/A}$ (0, 20.4.13), en les $d_{B/A}(s_i)$ et est par suite \emph{surjectif}; si $Y'=\Spec(B')$, et si $u:Y\to Y'$ est le $S$-morphisme correspondant à $\varphi$, on conclut de (17.2.2) que $u$ est \emph{non ramifié}. Si l'on prouve que le morphisme composé $u\circ g:X\to Y'$ est lisse au point $x$, il résultera donc de (17.7.10) que $u$ est étale au point $y$, puis de (17.3.5) que $g$ est lisse au point $x$; enfin, comme le morphisme structural $f':Y'\to S$ est lisse (17.3.8), $f=f'\circ u$ sera lisse au point $y$. Comme les images canoniques des $d_{B'/A}(T_i)$ dans $\Omega^1_{C/A}$ sont celles des $d_{B/A}(s_i)$, elles font partie d'une base de $\Omega^1_{C/A}$. On voit ainsi que pour prouver que c) implique a), on peut se borner au cas où $Y'=Y$, et par suite supposer que $f$ est un morphisme \emph{lisse}.

En vertu de (17.8.2), on peut alors se borner au cas où $S=\Spec(k)$ est le spectre d'un corps; de plus, grâce à (17.7.1, (ii)), on peut supposer que $k$ est algébriquement clos; enfin, remplaçant au besoin $X$ et $Y$ par des voisinages ouverts de $x$ et $y$ respectivement, on peut supposer que $f$ et $h$ sont tous deux lisses, et que l'homomorphisme canonique
\[
g^*(\Omega^1_{Y/S})\to\Omega^1_{X/S}
\]
est inversible à gauche (0, 19.1.12). Alors l'ensemble des points de $X$ rationnels sur $k$ est très dense dans $X$ (10.4.8), donc, pour prouver que $g$ est lisse, il suffit de prouver que $g$ est lisse en tout point $x\in X$ rationnel sur $k$, ou encore qu'en un tel point, $\sh{O}_{X,x}$ est un $\sh{O}_{Y,y}$-module plat et $\sh{O}_{X,x}/\mathfrak m_y\sh{O}_{X,x}$ un anneau régulier (17.5.1). Or, $y=g(x)$ est

